\documentclass[12pt]{article}
\usepackage[export]{adjustbox}
\usepackage{amsfonts, amsmath, amssymb}
\usepackage[english, spanish]{babel}
\usepackage{bm}
\usepackage{breqn}
\usepackage{enumitem}
\usepackage{fancyhdr}
\usepackage{geometry}
\usepackage{graphicx}
\usepackage{kantlipsum}
\usepackage{lastpage}
\usepackage{linebreaker} % line-breaker algorithm in LuaLaTeX
\usepackage{lualatex-math} % Fixes for mathematics
\usepackage{mathtools}
\usepackage{microtype}
\usepackage{nicefrac}
\usepackage[hyperref]{ntheorem}
% \usepackage{phfqit} % Bra-ket notation
\usepackage{physics}
\usepackage[nolabel]{showlabels}
\usepackage{siunitx}
\usepackage{titling}
\usepackage{tabularray}
\usepackage{tasks}
\usepackage{xcolor}
\usepackage{xfrac}
\usepackage{xparse}
\usepackage{xspace}
\usepackage{hyperref}
\usepackage{cleveref}

%%%%%%%%%%%%%%%%%%%%%%%%%%%%%%%%%%%%%%%%%%%%%%%%%%%%%%%%%%
%%%%%%%%%%%%%%%%%%%%%%%%%%%%%%%%%%%%%%%%%%%%%%%%%%%%%%%%%%
%%%%%%%%%%%%%%%%%%Configuraciones extras%%%%%%%%%%%%%%%%%%
\definecolor{base3}{RGB}{253, 246, 227}%
\definecolor{pinkwave}{RGB}{255, 0, 128}%
\definecolor{customBlue}{RGB}{11, 61, 98}%
\pagecolor{base3}
\graphicspath{{img/}}
\setlength{\parindent}{2em} % Sangría
\setlength{\parskip}{0.5em} % Espacio entre párrafos
\linespread{1.1} % line spacing
\setlength{\jot}{10pt} % Space between lines in multiline eqs
\crefname{equation}{ec.}{ecs.} % Equation's cross-reference name

% Line-breaker config
\linebreakersetup {
    maxtolerance = 90,
    maxemergencystretch = 1em,
    maxcycles = 4
}

%%%%%%%%%%%%%%%%%%Theorem environments%%%%%%%%%%%%%%%%%%
% Configuración de ambiente para problema
\theoremstyle{break}
\theoremheaderfont{\Large\normalfont\bfseries}
\theorembodyfont{\normalfont}
\theoremseparator{\bigskip} % Spacing between header and body
\theorempreskip{1.5em}
\theorempostskip{\topsep\bigskip}
\theorempostwork{
    \color{customBlue} \hrule width \hsize height 2pt \kern 1mm \hrule width \hsize
    }
\newtheorem{exercise}{Problema}
% Configuración de ambiente para solución
\theoremstyle{nonumberbreak}
\theoremheaderfont{\Large\normalfont\bfseries}
\theorembodyfont{\normalfont}
\theoremseparator{\medskip}
\theorempreskip{1em}
\theorempostskip{\topsep\medskip}
\newtheorem{solution}{Solución}

%%%%%%%%%%%%%%%%%%%%%%%%%%%%%%%%%%%%%%%%%%%%%%%%%%%%%%%%


% Configuración del paquete hyperref
\hypersetup{
    colorlinks = true,%
    linkcolor={[rgb]{0,0.2,0.6}},%
    citecolor={[rgb]{0,0.6,0.2}},%
    filecolor={[rgb]{0.8,0,0.8}},%
    urlcolor={[rgb]{0.8,0,0.8}},%
    runcolor={[rgb]{0.8,0,0.8}},% 
    menucolor={[rgb]{0,0.2,0.6}},%
    linkbordercolor={[rgb]{0,0.2,0.6}},%
    citebordercolor={[rgb]{0,0.6,0.2}},%
    filebordercolor={[rgb]{0.8,0,0.8}},%
    urlbordercolor={[rgb]{0.8,0,0.8}},%
    runbordercolor={[rgb]{0.8,0,0.8}},%
    menubordercolor={[rgb]{0,0.2,0.6}},% 
    pdftitle={Tarea X},%
    pdfauthor={López Merino Marcos},%
    pdfsubject={Subject},%
    pdfkeywords={Facultad de Ciencias, UNAM, materia, palabras clave},%
    unicode = true%
}

%%%%%%%%%%%%%%%%%%siunitx configuration%%%%%%%%%%%%%%%%%%
% Configuración del paquete siunitx
\sisetup{
	output-decimal-marker = {.}, 
	per-mode = symbol-or-fraction,
	separate-uncertainty = false,
	exponent-product = \cross,
    % inter-unit-product = \ensuremath{{}\vdot{}}
}

% Declaring new units
\DeclareSIUnit\kilogram{\kilo\gram}

%%%%%%%%%%%%%%%%%%%%%%%%%%%%%%%%%%%%%%%%%%%%%%%%%%%%%%%%

% Geometría del documento
\geometry{
    letterpaper,
    top = 0.6in,
    bottom = 0.8in,
    left = 0.6in,
    right = 0.6in,
    footskip = 38pt
}

%%%%%%%%%%%%%%%%%%Nuevos comandos%%%%%%%%%%%%%%%%%%
\newcommand*{\group}{8231}
\newcommand*{\classname}{Relatividad}
\newcommand*{\homeworknumber}{Tarea 1}
\newcommand*{\name}{Marcos López Merino}

% unit vector i
\newcommand*{\uveci}{{\bm{\hat{\textnormal{\bfseries\i}}}}}
% unit vector j
\newcommand*{\uvecj}{{\bm{\hat{\textnormal{\bfseries\j}}}}}
% unit vector
\DeclareRobustCommand{\uvec}[1]{{%
  \ifcsname uvec#1\endcsname
     \csname uvec#1\endcsname
   \else
    \bm{\hat{\mathbf{#1}}}%
   \fi
}}% 
\newcommand{\idest}{\emph{i.e.},\xspace} % id est
% Espacio vectorial, e.g., ℝ, ℂ, ℕ, etc.
\NewDocumentCommand{\vecspace}{m o}{%
  \IfValueTF{#2}{%
    \mathbb{#1}^{#2}%
  }{%
    \mathbb{#1}%
  }%
}
\newcommand*{\e}{\mathrm{e}} % exponential
% \newcommand{\solution}{\vspace*{4pt} \textbf{Solución}\vspace*{4pt}}

%%%%%%%%%%%%%%%%%%Portada y configuración%%%%%%%%%%%%%%%%%%
% Configuración de portada
\setlength{\droptitle}{-60pt} % raise the title

% Portada
\title{
    \textbf{\homeworknumber}\\
    \normalsize\vspace{0.1in}\small{\textbf{Entrega}:~\today}
    \vspace{-1.5in}
}
\author{}
\date{}

%%%%%%%%%%%%%%%%%%Header and footer%%%%%%%%%%%%%%%%%%
\setlength{\headheight}{15.2pt}
\pagestyle{fancy}
\lhead{Grupo \group, Sem. 2023-1}
\chead{\classname}
\rhead{\name}
\lfoot{\includegraphics[scale = 0.2, valign = c]{LogoFCUNAMcolor.pdf}}
% \lfoot{\includegraphics[scale = 0.1, valign = c]{example-image}}
\cfoot{\homeworknumber}
\rfoot{Pág. \thepage \hspace{1pt} de \pageref{LastPage}}

\renewcommand{\headrulewidth}{0.5pt}
\renewcommand{\footrulewidth}{0.5pt}
%%%%%%%%%%%%%%%%%%%%%%%%%%%%%%%%%%%%%%%%%%%%%%%%%%%%%%%%%%
%%%%%%%%%%%%%%%%%%%%%%%%%%%%%%%%%%%%%%%%%%%%%%%%%%%%%%%%%%

\begin{document}
    \maketitle
    \thispagestyle{fancy}
    
    
    \begin{exercise}
        Durante el curso usaremos unidades naturales, en las que \(c = 1\) y medimos tiempo y distancia en metros. Para que te acostumbres a esto en este problema harás algunas conversiones entre unidades naturales y el SI. Aunque parezca a primera vista un simple ejercicio mecánico, recuerda que \(c = 1\) vino de entender que el espacio y el tiempo son la misma cosas física, así que posiblemente escribir otras cantidades físicas en unidades naturales nos revele más de lo que pensamos.

        \begin{enumerate}[label = \alph*)]
            \item Transforma las siguientes cantidades del SI a unidades naturales. Expresas tus resultados en términos de kg y m.
            
            \begin{enumerate}[label = \arabic*.]
                \item Energía \(E = \SI{5}{\joule}\) (este es particularmente interesante, aquí hay implícito otro cambio total de paradigma).
                \item Momento \(p = \SI[per-mode = power]{3e4}{\kilogram \metre \per\second}\) (compara con tu resultado anterior, qué curioso).
                \item Densidad de masa \(\rho = \SI[per-mode = power]{10}{\kilogram \per \cubic\metre}\)
            \end{enumerate}

            \item Transforma las siguientes cantidades de unidades naturales al SI.
            
            \begin{enumerate}[label = \arabic*.]
                \item Velocidad \(v = \num{e-2}\) (siempre hablaremos de velocidades entre 0 y 1, así que es importante imaginarse cuál es la magnitud de estas en unidades más humanas).
                \item Presión \(P = \SI[per-mode = power]{e19}{\kilogram \per \cubic\metre}\)
                \item Densidad de energía \(U = \SI[per-mode = power]{1}{\kilogram \per \cubic\metre}\) (las mismas unidades que la presión, qué raro).
            \end{enumerate}

            \item Dos de las ecuaciones de Maxwell relacionan los campos eléctrico y magnético. En el vacío estas son
            
            \begin{dmath}
                \curl{\va*{E}} = -\pdv{\va*{B}}{t} \condition*{\curl{\va*{B}} = \dfrac{1}{c^{2}} \displaystyle\pdv{\va*{E}}{t}}
            \end{dmath}

            ¿Cómo se ven estas ecuaciones escritas en unidades naturales? ¿Cuál es la relación entre las unidades de \(\va*{E}\) y las de \(\va*{B}\)?
        \end{enumerate}
    \end{exercise}

    \begin{exercise}
        \kant[1]
        \begin{solution}
            \kant[1]
        \end{solution}
    \end{exercise}

    \begin{exercise}
        \kant[1]
        \begin{solution}
            \kant[1-3]
        \end{solution}
    \end{exercise}

    \pagebreak
    \begin{exercise}
        \kant[1]
        \begin{solution}
            \kant[1-4]
        \end{solution}
    \end{exercise}
\end{document}